\subsection{交换群的情形}

设 G 为局部紧群，\(N \subset \Gamma^{0}\) 为 G 的离散子群上的规范化 Haar 测度所成的子空间，\(N_{c}\) 为 N 中对应于使 G/H 为紧的 G 的离散子群 H 的子集；于是在 No.~3 命题 6 的记号下有 \(N_{c} = N \cap \Gamma_{c}^{0}\) ；并且若群 G 由 e 的一个紧邻域生成，则由 No.~3 命题 6 推出，\(N_{c}\) 在 N 中是开的（从而由 No.~3 命题 7 的推论它是局部紧的），并且在 \(N_{c}\) 上映射 \(\alpha \mapsto \|\mu_{\alpha}\|\) 的限制是淡连续的。

命题 9. --- 设 G 为由 e 的一个紧邻域生成的局部紧交换群。为使 \(N_{c}\) 的子集 A 在 \(N_{c}\) 中相对紧，必须且只需它满足下列两个条件：

\begin{enumerate}
\def\labelenumi{(\roman{enumi})}
\item
  存在一个开邻域 \(\mathbf{U}\) （\(e\) 在 \(\mathbf{G}\) 中的），使得 \(\mathrm{H}_{\alpha} \cap \mathrm{U} = \{e\}\) 对所有 \(\alpha \in \mathbf{A}\) 都成立。
\item
  存在一个常数 \(k\) ，使得 \(\mu_{\alpha}(\mathrm{G} / \mathrm{H}_{\alpha}) \leqslant k\) 对所有 \(\alpha \in \mathbf{A}\) 都成立。若 \(\mathbf{A} \subset \mathbf{N}_c\) 在 \(\mathbf{N}_c\) 中相对紧，则它更在 \(\mathbf{N}\) 中相对紧，故条件 (i) 与 (ii) 的必要性由 No.~3 命题 6 与 7 得出（无需假设 \(\mathbf{G}\) 交换）。反之，设 \(\mathbf{A} \subset \mathbf{N}_c\) 满足这些条件；若 \(\overline{\mathbf{A}}\) 是 \(\mathbf{A}\) 在 \(\mathbf{N}\) 中的闭包，则由 No.~3 命题 7，\(\overline{\mathbf{A}}\) 是紧的；此外，由于 \(\alpha \mapsto \| \mu_{\alpha} \|\) 在 \(\Gamma^0\) 上关于淡拓扑是下半连续的 (No.~2, Prop. 4)，条件 (ii) 蕴含 \(\| \mu_{\alpha} \| \leqslant k\) 也对所有 \(\alpha \in \overline{\mathbf{A}}\) 成立。而由于 \(\mathbf{G}\) 交换，\(\mu_{\alpha} = \mu / \alpha\) 是群 \(G / H_{\alpha}\) 上的一个 Haar 测度，因此 \(G / H_{\alpha}\) 对每个 \(\alpha \in \overline{\mathbf{A}}\) 都是紧的 (Ch. VII, §1, No.~2, Prop. 2)。这就表明 \(\overline{\mathbf{A}} \subset N_c\) ，从而 \(\mathbf{A}\) 在 \(N_c\) 中相对紧。
\end{enumerate}

推论. --- 设 G 为由 e 的一个紧邻域生成且满足 No.~4 的条件 (L) 的局部紧交换群。

设 \(D_c\) 为使 G/H 为紧的 G 的离散子群 H 所成的集合，并且对每个 \(H \in D_c\) ，设 \(v(H)\) 为总质量 \(\mu_\alpha(G/H)\) ，其中 \(\mu_\alpha\) 是 \(\mu\) 关于 H 的规范化 Haar 测度 \(\alpha\) 的商测度。为使空间 \(D_c\) 的子集 A 在 \(D_c\) 中相对紧，必须且只需它满足下列两个条件：

\begin{enumerate}
\def\labelenumi{(\roman{enumi})}
\item
  存在 e 在 G 中的一个开邻域 U，使得 \(\mathbf{H} \cap \mathbf{U} = \{e\}\) 对所有 \(\mathbf{H} \in \mathbf{A}\) 都成立。
\item
  存在一个常数 \(k\) ，使得 \(v(\mathbf{H}) \leqslant k\) 对所有 \(\mathbf{H} \in \mathbf{A}\) 都成立。
\end{enumerate}

考虑到命题 9，这立即由下述事实得出：\(D_{c}\) 是 \(N_{c}\) 在 N 到 D 的典范双射下的像，并且在所作假设下该双射是一个同胚 (No.~4, Prop. 8)。

例. --- 取 \(G = R^{n}\) ，并取 \(\mu\) 为 Lebesgue 测度；命题 9 的推论的全部假设都得到满足。使 G/H 为紧的 G 的离散子群 H 正是秩为 n 的离散子群 (GT, VII, §1, No.~1, Th. 1)；这样的子群 H 由一个基 \((a_{i})_{1 \leqslant i \leqslant n}\) （\(R^{n}\) 的基）生成，并且

\[
v (\mathbf {H}) = | \det (a _ {1}, \dots , a _ {n}) |
\]

（行列式是关于 \(\mathbf{R}^n\) 的标准基取的）(Ch. VII, §2, No.~10, Th. 4)。空间 \(\mathbf{D}_c\) 在这里可按如下方式解释：每个子群 \(\mathbf{H} \in \mathbf{D}_c\) 都是 \(g \cdot \mathbf{Z}^n\) 型的子群（即子群 \(\mathbf{Z}^n\) 经某个元素 \(g \in \mathbf{GL}(n, \mathbf{R})\) 变换而得），而 \(\mathbf{GL}(n, \mathbf{R})\) 中使 \(\mathbf{Z}^n\) 保持稳定的子群可与 \(\mathbf{GL}(n, \mathbf{Z})\) 等同。因此，\(\mathbf{D}_c\) 作为（非拓扑的）齐性空间可以典范地与 \(\mathbf{GL}(n, \mathbf{R}) / \mathbf{GL}(n, \mathbf{Z})\) 等同。另一方面，\(\mathbf{GL}(n, \mathbf{R})\) 连续地作用在 \(\mathbf{R}^n\) 中，从而也作用于关于淡拓扑的 \(\mathcal{M}_{+}(\mathbf{R}^{n})\) 中 (§3, No.~3, Prop. 13)，因而作用于子空间 \(\mathrm{N}_c\) （\(\mathcal{M}_{+}(\mathbf{R}^{n})\) 的）中；此外，\(\mathrm{N}_c\) 到 \(\mathrm{D}_c\) 上的典范同胚 (No.~4, Prop. 8) 与 \(\mathbf{GL}(n, \mathbf{R})\) 的运算律相容。由于 \(\mathbf{GL}(n, \mathbf{R})\) 在无穷处可数且 \(\mathrm{D}_c\) 局部紧，上面定义的 \(\mathbf{GL}(n, \mathbf{R}) / \mathbf{GL}(n, \mathbf{Z})\) 到 \(\mathrm{D}_c\) 上的双射是一个同胚 (Ch. VII, App. I, Lemma 2)。于是命题 9 的推论给出了齐性空间 \(\mathbf{GL}(n, \mathbf{R}) / \mathbf{GL}(n, \mathbf{Z})\) 中紧性的一个判别法。

